\chapter{Heat-Kernel Response to an External Field Strength}
\label{app:heat_kernel_mixing}

This appendix records the leading response of the Seeley--DeWitt
coefficient \(a_1\) to a laboratory field strength and the precision
constraints that any residual conformal mixing must satisfy. It does
not recompute \(a_1\) on the locked block, does not derive a kinetic
mixing from the entanglement functional, and does not open an
observationally allowed laboratory window.

\section{Vanishing of the linear term in \(a_1\)}
\label{sec:a1-null}

On a Laplace-type operator \(\Delta=-D^\mu D_\mu+E\) the first
Seeley--DeWitt coefficient is
\begin{equation}
  a_1=\operatorname{tr}(E)-\frac16 R\,\dim.
\end{equation}
If a controllable field strength enters the endomorphism only through
an antisymmetric generator,
\begin{equation}
  E=E_0+\xi F_{\alpha\beta}S^{\alpha\beta},
\end{equation}
the trace of \(S^{\alpha\beta}\) vanishes on every representation that
appears in the locked \(4\times 4\) block and in the diffeomorphism
and local-Lorentz ghost sectors. Consequently
\begin{equation}
  \delta a_1=0.
\end{equation}
This is a standard algebraic fact. It is not a recomputation of
\(a_1\) on the NAMCG fluctuation operator (Problem~F remains open),
nor does it cancel the higher-order suppression by
\(\Lambda_{\mathrm{UV}}\) already identified for laboratory fields.

\section{Quadratic appearance in \(a_2\)}
\label{sec:a2}

The first non-vanishing correction linear in the laboratory field
appears in \(a_2\) and is proportional to \(F_{\mu\nu}F^{\mu\nu}\).
Variation of the resulting effective action with respect to the
metric produces a contribution proportional to \(g_{\mu\nu}\). The
coupling is therefore confined to the conformal mode. No
transverse-traceless response is generated at this order.

\section{Precision constraints on conformal mixing}
\label{sec:constraints}

Any additional kinetic mixing between the conformal mode \(\psi\) and
a laboratory scalar that might be introduced by hand is constrained
by existing data:
\begin{itemize}
  \item Cassini: \(|\gamma-1|\lesssim 2\times 10^{-5}\);
  \item torsion-balance equivalence-principle tests:
        \(\sin^2\theta\lesssim 10^{-13}\).
\end{itemize}
A mixing angle large enough to produce a macroscopic laboratory
acceleration violates these bounds. No calculation within the
present framework cancels the parametric suppression
\(F/\Lambda_{\mathrm{UV}}^2\) while leaving a residual that survives
the bounds. The observational window therefore remains empty.

\section{What is not shown}
\label{sec:not-shown}

\begin{itemize}
  \item A completed evaluation of \(a_1\) (or the relevant part of
        \(a_2\)) on the locked \(4\times 4\) block with ghosts.
  \item A derivation of a kinetic mixing term from
        \(U_{\mathrm{ent}}[I(A:B)]\) or from the modular Hamiltonian.
  \item A dynamical cancellation of quadratic ultraviolet divergences
        that leaves a finite infrared residual.
  \item Any projection of galactic rotation-curve discrepancies onto
        solar-system scales, or any fit to pulsar-timing-array
        correlation data. Those exercises do not follow from the
        heat-kernel expansion and are omitted.
\end{itemize}

\section{Status}
\label{sec:appK-status}

The linear coupling of a laboratory field strength to \(a_1\)
vanishes for antisymmetric generators, and any residual conformal
mixing is limited by existing precision tests to a level far below
laboratory interest. This appendix adds no new handle on the
entanglement stress or on \(G_{\mathrm{eff}}\) and does not alter
the open-problem ledger of Appendix~\ref{app:status}.